\begin{proposition}[7.5.6]
\label{IV.2.s7.5.6-fr}
Sean $k$ un cuerpo y $A$, $B$ dos anillos locales noetherianos completos
que contienen $k$ y cuyos cuerpos residuales son extensiones finitas
de $k$. Sea, por otra parte, $\boldsymbol R$ una propiedad que satisface
las condiciones $(\mathrm R'_{\mathrm I})$,
$(\mathrm R_{\mathrm{III}})$ y $(\mathrm R_{\mathrm{IV}})$ \textup{(la propiedad
$\boldsymbol P$ se define a partir de $\boldsymbol R$ por (7.5.0))}.
Supongamos que se cumplen $\boldsymbol R(A)$ y
$\boldsymbol P(\operatorname{Spec}(B),k)$ \textup{(lo que significa que, para
toda extensión finita $k'$ de $k$ y para cada uno de los anillos locales
completos $B'_i$ cuyo producto directo es $B\otimes_k k'$, se cumple
$\boldsymbol R(B'_i)$)}. Entonces, para cada uno de los anillos locales
completos $C_j$ cuyo producto directo es $A\widehat\otimes_k B$,
se cumple $\boldsymbol R(C_j)$.
\end{proposition}

\begin{proof}
De (7.5.5, (iii)) se deduce que cada uno de los homomorfismos
$A\to C_j$ es local, y de (7.5.5, (ii)), que cada uno de los $C_j$ es
un $A$-módulo plano; por último, por (7.5.5, (iv)), los cuerpos
residuales de los $C_j$ son extensiones finitas de $A/\mathfrak m$.
Por otra parte, la propiedad
$\boldsymbol P(\operatorname{Spec}(C\otimes_A(A/\mathfrak m)),A/\mathfrak m)$
es verdadera, pues $C/\mathfrak mC=(A/\mathfrak m)\otimes_{\boldsymbol{k}}B$, y la
afirmación se deduce de la definición de $\boldsymbol P$ y de la
hipótesis de que $A/\mathfrak m$ es una extensión \emph{finita} de $k$.
Se satisfacen, por tanto, todas las condiciones de (7.5.1) para los
homomorfismos locales $A\to C_j$, y los morfismos correspondientes
$\operatorname{Spec}(C_j)\to\operatorname{Spec}(A)$ son, por tanto,
$\boldsymbol P$-\emph{morfismos}; la conclusión se deduce entonces de
$(\mathrm R'_{\mathrm I})$ y de la hipótesis de que se cumple
$\boldsymbol R(A)$.
\end{proof}
